\section{Protobuf Transformation}\label{sec:proto-transform}

At the core of the Protobuf compatibility methodology is a transformation which
projects the view of a message under one descriptor to a related message under
another descriptor.

\subsection{Origins}

Using a function to transform the value isn't an arbitrary technical choice, nor
is it only useful to model potential descriptor changes to specific Protobuf
value. It arises naturally from the proofs of the intermediate format, and that
evolution is described here.

The first and strongest intermediate parser compatibility theorem is the
schema-correct one, Theorem~\ref{thm:inter-schema-correct}, where the relations
at play are strong enough that we can directly prove that the resulting value
from a serializer / parser round-trip under the same descriptor is equal. No
transformation is needed here.

However, the next step was to weak the compatibility antecedent to arrive at
Theorem~\ref{thm:inter-id-compat} where the descriptor is still the same but can
have no relation to the message being operated on. This is no longer possible
with our new totality conditions (Section~\ref{sec:value-totality}), but led to
the first example of a transformation function. The Intermediate ID
Compatibility relation has the \textsc{M-Drop-Unknown} rule, with no constrain
which value is dropped, and \textsc{M-Missing}, which can insert a new value
into the resulting message. Because of these rules, the induction hypothesis of
the theorem by itself is too weak to prove the whole theorem since the
compatibility relation has multiple $m_2$ related to the same $m_1$, $d_1$ and
$d_2$ so the proof must find a witness. Instead we use the transformation to
strengthen the inductive hypothesis into naming that witness and prove that
anything from the transformation will land inside the compatibility relation.

The final intermediate format theorem, Theorem~\ref{thm:inter-compat},
reinforces this. It uses the same technique with a much more complex
transformation function that actually mirrors the Protobuf transformation. Both
are a structural recursion on the reader's descriptor field list, building the
new interpretation of the value under the new descriptor using the writer's
descriptor. The Protobuf one also has the stronger justification of enforcing
the totality requirement.

\subsection{Definition}

The full definition of the transformation function is omitted, since it is not
fully defined at the time of this writing. However, it is broken down into
levels, just like the various size functions.

\begin{minted}{lean}
Value.reinterpret        (d₁ d₂ : Desc) (v : Value) : Value
Value.reinterpretEntries (d₁ : Desc) (v : Value) :
    List ((_:Int) × Field) → List ((_:Int) × Val)
Value.reinterpretAt      (d₁ : Desc) (v : Value) (k : Int) (f₂ : Field) : Val
Slot.reinterpret         (f₁ f₂ : Field) (x : Val) : Option Val
Payload.reinterpret      (t₁ t₂ : FieldType) (p : Payload) : Option Payload
\end{minted}

The \texttt{reinterpretEntries} function is basically a \texttt{map} over the
internal association list, with the required termination metric.

The \texttt{Option} in the last two function signatures is \emph{not} a
message-level optional; returning a \texttt{none} means that the value is not
carryable by the new type and the caller falls back to the default slot value,
which is exactly what the reader sees when the writer's bytes do not populate
the field at all or using the incorrect wire type.

Mathematically, the transformation where a message $m$ at descriptor $d_1$ is
transformed to $d_2$ is denoted $\transform{d_1}{m}{d_2}$.

For each key in the reading descriptor, there are three cases.

\begin{itemize}
\item \textbf{Shared.} If both descriptors carry the payload, move it over to
  the new message, transforming it if needed.
\item \textbf{Reader only.} Inject the default value for that field.
\item \textbf{Writer only.} This is a dropped field, and since we transform the
  message by walking the \emph{reader's} fields, this case cannot arise.
\end{itemize}

The other possible cases are that there is a value in the message which isn't in
the writing descriptor or that there is a field in the writing descriptor which
isn't in the message. The totality assumption from
Section~\ref{sec:value-totality} eliminates these cases from valid messages.

The recursion here is driven by the reader's descriptor since we want to produce
a message that's total over the reader's fields. Thus the termination metric is
the size that descriptor alone, not dependent on the message being
reinterpreted.

Note that the intermediate compatibility relation didn't have a drop rule in
$\ll$. That was because if we encountered a dropped value, we had no way to know
if the next thing in the encoding was a scalar or a nested message and the
parsing failed. Since Protobuf encoding the wire type into the value encoding,
it's always possible to skip the unknown field without de-synchronizing the byte
stream. The Protobuf drop rule will be scoped to only be applied only when
$\dom{d_1} \not \subset \dom{d_2}$.

One final note is that the measure chain used by all the transformation layers
is only \emph{weakly} decreasing within one layer, so the measures are scaled by
4 with a per-step offset to ensure that the termination metric will actually
decrease with each application. 

\vspace{1em}
\begin{theorem}[Transformation Specification]
  For all messages $m$ and descriptors $d_1$, $d_2$, each value of the
  transformed message is equal to the result of transforming that value.
  \[ \transform{d_1}{m}{d_2}[k] = \transform{d_1}{m[k]}{d_2} \]
\end{theorem}

\vspace{1em}
\begin{theorem}[Transform WF]\label{thm:trans-wf}
  \[ \forall\ m, d_1, d_2.\ d_2 \checkmark_w \rightarrow \transform{d_1}{m}{d_2} \checkmark_w \]
\end{theorem}

\vspace{1em}
\begin{theorem}[Transform Total]\label{thm:trans-total}
  \[ \forall\ m, d_1, d_2.\ \dom{d_2} = \dom{\transform{d_1}{m}{d_2}} \]
\end{theorem}

\vspace{1em}
\begin{theorem}[Transform Self]\label{thm:trans-self}
  \[ \forall\ m, d.\ d \checkmark \rightarrow m \checkmark^d \rightarrow \transform{d}{m}{d} = m \]
\end{theorem}

\vspace{1em}
\begin{theorem}[Transform Valid]\label{thm:trans-valid}
  \[ \forall\ m, d_1, d_2.\ m \checkmark^{d_1} \rightarrow d_2 \checkmark \rightarrow \mathrm{OneofPreservedAll}(d_1, d_2) \rightarrow \transform{d_1}{m}{d_2} \checkmark^{d_2} \] 
  Where
  \begin{itemize}
  \item $\mathrm{OneofPreservedAll}(d_1, d_2) = \mathrm{OneofPreserved}(d_1, d_2) \wedge
    \forall\ k, c_1, d_1', c_2, d_2'.\ d_1[k] = \langle c_1, d_1' \rangle \rightarrow d_2[k] = \langle c_2, d_2' \rangle \rightarrow
    \mathrm{OneofPreservedAll}(d_1', d_2')$
  \item $\mathrm{OneofPreserved}(d_1, d_2) = \forall\ k_1, k_2, f_1, f_1', f_2, f_2',
    g. d_1[k_1] = f_1 \rightarrow d_1[k_2] = f_2 \rightarrow d_2[k_1] = f_1' \rightarrow d_2[k_2] = f_2' \rightarrow
    f_1' = \langle \oneof{g}, \_ \rangle \rightarrow f_2' = \langle \oneof{g}, \_ \rangle \rightarrow \exists\ g_0. f_1 = \langle
    \oneof{g_0}, \_ \rangle \wedge f_2 = \langle \oneof{g_0}, \_ \rangle$
  \end{itemize}
\end{theorem}

The two important theorems here are Theorem~\ref{thm:trans-self}, which states
that transforming a message against it's own transform produces the same
message, and Theorem~\ref{thm:trans-valid} which we can use to show that a
transformed value is valid.

Most of the premises for these theorems should be evident, except the
\texttt{oneof} conditions for Theorem~\ref{thm:trans-valid}. This premise is
required, and to see why consider that a writer should legitimately set two
fields outside a \texttt{oneof} and then transform to a reader where those two
fields are part of the same \texttt{oneof}. The OneofPreserved function states
that any two reader keys sharing a \texttt{oneof} group were already in a group
together in the writer. Note that this is vacuous if either key is reader-only,
so adding a new field to an existing \texttt{oneof} is acceptable, while moving
an existing field into an existing \texttt{oneof} isn't. This is part of the
Protobuf best practices guide~\cite{LanguageGuideProto}, represented as proof
obligation here. The direction is also important, since it does allow us to
split a \texttt{oneof} into multiple \texttt{oneof} groups without
issue. Additionally, the recursive version is required to prevent the example
case from appearing in a nested message.

Now, it is possible to move an existing field into an existing \texttt{oneof} in
a Protobuf descriptor. In that case, you'd need to be sure that the value didn't
set multiple fields in a \texttt{oneof} of the reader (we'll know in advance
that it can't have multiple fields set in a writer's \texttt{oneof}). Then you
might be able to weaken the $\mathrm{OneofPreservedAll}(d_1, d_2)$ to 
\[ \forall\ k_1, k_2, g.\ d_2[k_1] = \langle \oneof{g}, \_ \rangle \rightarrow d_2[k_2] = \langle \oneof{g}, \_ \rangle \to
  m[k_1] = \mathtt{NONE} \vee m[k_2] = \mathtt{NONE} \]
Which is wouldn't be true for all messages under $d_1$, but might be provable
as a weakened version of Theorem~\ref{thm:trans-valid} for some systems.

{\color{red} This is another place were the functional transformation breaks
  down. In the lean serializer, everything is serialized in sorted field key
  order, so the last writer wins rule (Section~\ref{sec:last-writer-wins}) would
  be deterministic. However, since a real serializer could serialize in an
  arbitrary order, this is actually a non-deterministic operation. The
  relational encodes definition would allow this and further complicate the
  non-canonical Protobuf encoding.}

\subsection{UPDATE:\@ What's implemented right now}

At the moment, the transformation supports width and depth changes to messages
(i.e.\ adding new fields or changing nested messages) and moving fields between
\texttt{optional} and \texttt{oneof}. Scalar retyping is not implemented yet,
but will be.

The one supported cardinality change, moving between an \texttt{optional} and
\texttt{oneof} naturally requires that we enforce the \texttt{oneof} requirement
that only one field in the \texttt{oneof} is set.

\subsection{Faithfulness Theorem}

The Protobuf compatibility framework might drop the message compatibility
relation in favor of only having the transformation. In that case, we'd end up
needing to argue that the transformation is total over all possible messages
being converted between two descriptors, i.e. $\transform{d_1}{m}{d_2}$ produces
some $m'$ for all $m \checkmark^{d_1}$.

The current transformation function is \emph{always total} by falling back to
injecting field type defaults, so in the context of Protobuf, we want to show
that we can produce a total transformation without needing this behavior outside
of injecting the new values for new fields (since any other place were a default
value is injected would imply that something failed to reinterpret). Note that
this theorem can't be solely stated about the output of
$\transform{d_1}{m}{d_2}$ since the output can't differentiate between the
default value naturally occurring and being injected due to a problematic
update.

\begin{minted}{lean}
theorem Slot.reinterpret_isSome {f₁ f₂ : Field} {x : Slot}
    (h : f₁ ∝ f₂) (hm : Slot.Matches x f₁) :
    (Slot.reinterpret f₁ f₂ x).isSome
\end{minted}

Including \texttt{Slot.Matches} is important since incompatible cardinalities or
types between the slot and the original descriptor would invoke the default
value fallback anyways. This could be built into a faithfulness theorem.

\vspace{1em}
\begin{theorem}[Faithfulness]\label{thm:faith}
  For two compatible descriptors, $d_1 \ll d_2$ and a valid value $m$ under
  $d_1$, $m \checkmark^{d_1}$ such that for all keys $k$, $d_1[k] = f_1$ and $d_2[k] =
  f_2$ there must exist a value $x$ and $y$ such that
  \begin{itemize}
  \item $m[k] = x$
  \item $\transform{f_1}{m}{f_2} = y$
  \item $\transform{d_1}{m}{d_2}[k] = y$
  \end{itemize}
\end{theorem}

Which basically characterizes the reinterpretation process as point-wise over the
message.

\begin{minted}{lean}
theorem Value.get?_reinterpret_of_compat {d₁ d₂ : Desc} {v : Value} {k f₁ f₂}
    (h : d₁ ≪ d₂) (hv : Value.Valid d₁ v)
    (h₁ : d₁.get? k = some f₁) (h₂ : d₂.get? k = some f₂) :
    ∃ x y, v.get? k = some x ∧ Slot.reinterpret f₁ f₂ x = some y ∧
      (Value.reinterpret d₁ d₂ v).get? k = some y
\end{minted}

Some corollaries from this include the fact that repeated slots don't drop
values. If the scalar transformation is only the identity function, which is
presently the case, we also have that
$\transform{d_2}{\transform{d_1}{m}{d_2}}{d_{2}} = m$ since the newly injected
fields are simply dropped in the backwards directly. This means that no data is
lost, but the statement is too strong to handle scalar conversion and in
particular integer width changes which might truncate values.

\vspace{1em}
\begin{theorem}[Point-wise Field Compatibility]\label{thm:pointwise-compat}
  \[ d_1 \checkmark \rightarrow \forall\ k \in \dom{d_1} \cap \dom{d_2}. d_1[k] \propto d_2[k] \] 
\end{theorem}

Combined with Theorem~\ref{thm:faith}, Theorem~\ref{thm:pointwise-compat} states
that the descriptor relation $\ll$ is a maximal relation or equivalently the
weakest precondition for lossless data transmission. Again, this will be
complicated by adding integer truncation. Once that feature is added, this
concept is split into \emph{carriage}, which thinks about data being
present, and \emph{lossless-ness} that the transformed view of the data is
unique. 

% Local Variables:
% citar-bibliography: ("../pollux.bib")
% TeX-master: "../pollux.tex"
% TeX-command-default: "Make"
% jinx-local-words: "protobuf"
% End:
